\chapter{Conclusions and Future Work}

\section{Conclusions}

This thesis set out to build a FHIR-native, knowledge-graph-based clinical recommendation
system for chronic coronary disease, and to establish by measurement what the FHIR-native
design actually buys. The system was built, and each of the three research objectives stated
in Chapter~I can now be answered from the results of Chapter~IV.

\begin{enumerate}
    \item \textbf{The pipeline was realised and evaluated end to end.} HL7 FHIR~R4 bundles
    for 5{,}490 patients were loaded into a Neo4j property graph and fused with the 2023
    AHA/ACC CCD guideline and a guideline-seeded literature corpus through a UMLS concept
    bridge, yielding a graph of 1{,}804{,}379 nodes and 8{,}796{,}425 relationships over 25
    labels, 51 relationship types and 26 code systems
    (Section~\ref{sec:kg-results}). The \texttt{:Patient} count equals the number of patients
    Synthea reported generating, so the idempotent load neither dropped nor duplicated a
    record. Against a reference standard generated over a data path separate from the system's,
    guideline concordance over the 876-patient evaluation cohort reached
    $F_1 = 0.810$ (Section~\ref{sec:ablation}); with the 15 recommendations whose predicate
    both sides share removed, the remaining 53 reach $F_1 = 0.797$
    (Section~\ref{sec:shared-predicate}). That figure rests on a
    validated reference rather than an assumed one: on the 30 patients labelled
    manually, the automated standard agrees on 99.5\% of cells
    (2{,}150 of 2{,}160) with Cohen's $\kappa = 0.991$ (95\% confidence interval
    0.984--0.996, resampling patients; Section~\ref{sec:kappa-ci}). Thirty patients are enough
    to support the claim of near-perfect agreement, but they were labelled by a single annotator
    (Section~\ref{sec:limitations}). The clinician review adds a judgement the reference
    standard cannot give: a cardiologist judged 77.7\% of 247 recommendations for 15 patients
    appropriate (Section~\ref{sec:clinician}). Outside the antithrombotic and beta-blocker
    recommendations, almost every recommendation was judged appropriate. The ratings of
    inappropriateness and the disagreement between raters ($\kappa = 0.29$) concentrate on
    antithrombotic recommendations, whose guideline logic depends on time since an acute
    coronary syndrome, bleeding risk, and the indication for anticoagulation. Objective~1 is
    therefore met.

    \item \textbf{Observation gating improves guideline concordance, and the ablation
    localises the improvement exactly.} Gating observation-dependent recommendations on the
    patient's measured values raises $F_1$ from 0.743 to
    0.810, a gain of $+0.067$ obtained through precision
    (0.612 to 0.684) with recall essentially unchanged
    (Section~\ref{sec:ablation}). The classification of Section~\ref{sec:disc-class} makes a
    falsifiable prediction about where that gain must live, and it holds exactly: 53 of the
    68 recommendations are identical in both arms down to the individual true-positive,
    false-positive and false-negative counts, and the whole of the difference originates in
    the 15 that are gated. A second and less obvious result is that gating is not one
    mechanism but two. For the five \emph{value-gated} recommendations it acts on precision, which rises from 0.194
    to 1.000 because they no longer fire for every cohort member. For the ten in the
    \emph{both} class it acts overwhelmingly on recall, which rises from 0.006 to 1.000.
    Their encoded condition targets alone rarely match, so it is the numeric predicate that
    identifies the patients they apply to (Section~\ref{sec:ablation-class}). Describing the ablation solely as a precision
    improvement, as the aggregate invites, is therefore wrong for two thirds of the affected
    recommendations. One qualification attaches to the ablation as a whole. The fifteen gated recommendations are
    precisely those whose predicate the system and the reference standard share. What the
    ablation measures is therefore implementation fidelity rather than the clinical validity
    of the thresholds (Section~\ref{sec:shared-predicate}). Objective~2 is met.

    \item \textbf{Citation-level faithfulness is mechanically decidable, but semantic
    faithfulness is not guaranteed.} The grounded briefing retrieves every fact
    deterministically from the graph, carries a recommendation identifier, class of
    recommendation and PMIDs with each one, and confines the language model to phrasing them.
    All 30 generated briefings passed the automated check: no fabricated recommendation or
    PMID, no omitted recommendation, and no anticoagulation over-claim
    (Section~\ref{sec:faithfulness}). The structural part of this result is that the model
    sits downstream of retrieval with no path back to the graph, so any recommendation outside
    what retrieval supplied must be fabricated and will be detected. The check of citable items
    is therefore decidable, which similarity-based retrieval-augmented generation cannot offer.
    Deterministic retrieval does not guarantee semantic faithfulness: the model can still
    phrase correct facts wrongly, and this check catches only one pattern of such error. That
    pattern was itself corrected after its output was seen. Objective~3 is met within those
    bounds.
\end{enumerate}

Beyond the three objectives, one finding emerged that was not anticipated at the outset and
that bears on any FHIR-based deployment. The two projections of the same synthetic record
are not comparable in richness. The FHIR export retains a categorical medication status but
omits stop dates, while the CSV export retains stop dates but carries no status, so neither
is a superset of the other. The chronic-NSAID rule depends on how long a drug was taken, and
consequently only two of its four arms can be computed from the graph
(Section~\ref{sec:chronicity}). The cost is measured rather than hypothetical: the
chronicity-gated safety variant reaches precision 0.909 at recall
0.526, missing exactly those patients whose use is chronic only through a
single long course (Section~\ref{sec:safety-detection}). A system that consumes FHIR alone
must anticipate this class of information loss wherever a guideline depends on exposure
duration.

That asymmetry runs in the other direction too, and it narrows a claim this work originally
made. The precision gain from observation gating was attributed to the FHIR-native design from
Chapter~I onwards, but the attribution was asserted rather than tested. When the same gate was
rebuilt reading only the tabular export, not one of the 13{,}140 decisions compared changed, and
the cohort selected identically (Section~\ref{sec:csv-gating}). The gain itself is real and
unchanged. Its source was misstated: what produces it is loading numeric observation values as
first-class data, and FHIR is the most direct way to obtain them rather than the only way. The
conclusion was reached by testing the work's own claim instead of defending it, and a negative
result of this kind is worth reporting precisely because attributions of this sort are
commonly asserted without test in comparable work.

\section{Future Work}

The following directions follow from the limitations set out in
Section~\ref{sec:limitations} and are ordered by how much each would change the results.

\begin{enumerate}
    \item \textbf{Widen the cohort gate to a disjunction of concepts.} The CCD gate accepts
    a single UMLS concept, whereas the guideline defines chronic coronary disease
    disjunctively. Sixty-seven patients---$7.1\%$ of the clinically defined
    cohort---whose only CCD evidence is a prior-infarction or prior-revascularization
    \emph{situation} code are therefore never matched against any recommendation, and in a
    deployment would receive no output at all (Section~\ref{sec:gate}). This is a small
    change to one query and is the single largest available gain in coverage.

    \item \textbf{Extend the clinician review and conduct a briefing reader study.} The
    clinician review here covers 15 patients for one cardiologist, and agreement between
    raters is only fair (Section~\ref{sec:clinician}). A follow-up review should have several
    cardiologists rate every patient. The review also points to improvements in the system
    itself. The briefing's conclusion should name the flagged drugs and the information the
    antithrombotic and beta-blocker recommendations rest on, such as time since an acute
    coronary syndrome, ejection fraction, and heart-failure status. Where the record cannot
    establish that information, as with bleeding risk, the conclusion should say so. Missing
    information of this kind is what most often prevented a judgement. Separately,
    faithfulness is a property of the briefing's relation to the record rather than of its
    value to a reader. A study on the model of Van~Veen \textit{et al.}~\cite{vanveen2023}, in
    which clinicians score summaries for completeness, correctness and conciseness with
    inter-rater agreement quantified by Cohen's $\kappa$ \cite{cohen1960}, would close that
    gap. It would also lift the faithfulness evaluation from an automated check of citable
    items to a human judgement of semantic faithfulness.

    \item \textbf{Repeat the evaluation on real records.} Synthea patients are generated
    from disease-progression models, so their comorbidity patterns and laboratory
    trajectories are internally consistent but not drawn from a real population. The results
    reported here establish that the pipeline is internally correct; they do not establish
    that it would achieve the same concordance on records that are noisier, sparser and
    accompanied by unstructured text. Real records also carry the cardiac symptoms and
    complaints that Synthea does not, so the six angina-dependent recommendations can only be
    tested there (Section~\ref{sec:limitations}).

    \item \textbf{Recover duration-dependent guideline logic.} Either a FHIR profile that
    carries dispense or dosage periods, or the pairing of the graph with an auxiliary
    source, would make all four arms of the chronic-use rule computable and would remove the
    recall ceiling measured in Section~\ref{sec:safety-detection}.

    \item \textbf{Give the safety tier a richer representation.} Population-gated matching is perfectly precise but recovers only the minority of
    contraindications whose population condition is structurally gateable. The misses are
    the negation-defined and conjunctive contraindications, which a flat schema cannot
    express (Section~\ref{sec:safety-detection}). Representing those explicitly, rather than
    choosing between over-flagging and under-recovering, is the way past the trade-off.

    \item \textbf{Port the design to a second guideline or disease.} Only the 2023 AHA/ACC
    CCD guideline is encoded, and only for chronic coronary disease. The concept-bridge
    design makes generalisation plausible because nothing in the query layer is specific to
    cardiology, but this work does not demonstrate it; encoding a second guideline and
    re-running the same evaluation would.
\end{enumerate}
